% ==============================================================================
\chapter*{Nota sobre este borrador}
\addcontentsline{toc}{chapter}{Nota sobre este borrador}
% ==============================================================================

Este documento es un \textbf{borrador parcial} de la memoria del Trabajo Fin de Grado, preparado para una primera revisión por parte del tutor antes de la entrega definitiva (22~de~mayo de 2026). Su objetivo es validar la estructura, el enfoque y el alcance del proyecto, no presentar el documento completo.

\section*{Estado actual de la memoria completa}

La memoria completa cuenta a fecha de hoy con aproximadamente 89~páginas y 6 capítulos cerrados, junto con 3 anexos. Para no abrumar la primera revisión, este borrador incluye únicamente:

\begin{itemize}[leftmargin=1.2cm, itemsep=2pt]
  \item El \textbf{índice completo} de la memoria final, para revisar la estructura propuesta.
  \item El \textbf{Capítulo~1 completo} (Introducción), que contiene la motivación, los objetivos formales y la declaración de uso de IA.
  \item Una \textbf{síntesis del Capítulo~3} (Análisis y diseño), con la estructura por secciones y las cifras agregadas en lugar de las cinco tablas completas de requisitos. Las tablas detalladas están redactadas y disponibles en el documento completo.
  \item El \textbf{Capítulo~6} (Conclusiones, limitaciones y trabajo futuro).
\end{itemize}

Los Capítulos 2, 4 y 5, así como los Anexos A, B y C, están redactados pero se omiten en este borrador para facilitar una revisión enfocada. Pueden adjuntarse en próximas tutorías a petición del tutor.


\section*{Lo que aún está pendiente de incorporar}

\begin{itemize}[leftmargin=1.2cm, itemsep=2pt]
  \item \textbf{Capturas de pantalla} de la aplicación (8 capturas previstas: mapa, formulario de reporte, detalle, panel admin, registro, manual, vista móvil, configuración 2FA).
  \item \textbf{Fichas de las 5 aplicaciones del estado del arte}, tras probar personalmente FixMyStreet, Línea Verde, iNaturalist, SeeClickFix y Epicollect5.
  \item \textbf{Resultados de las pruebas de usabilidad} con tres usuarios reales.
  \item \textbf{Portada oficial EPS}, en sustitución de la provisional.
  \item \textbf{Texto de los agradecimientos}.
\end{itemize}


\section*{Cuestiones para la próxima tutoría}

\begin{itemize}[leftmargin=1.2cm, itemsep=2pt]
  \item Confirmación del sistema de citas (APA o IEEE).
  \item Confirmación de la modalidad de defensa (A o B).
  \item Decisión sobre publicación en RUA.
  \item Cualquier observación sobre alcance, estructura o correcciones a aplicar antes de la entrega.
\end{itemize}

\noindent\rule{\textwidth}{0.4pt}
\begin{flushright}
\small Erardo Aldana Pessoa\\
Convocatoria C3 — Curso 2025--2026
\end{flushright}
